% !TEX TS-program = pdflatex
% !TEX encoding = UTF-8 Unicode

% This file is a template using the "beamer" package to create slides for a talk or presentation
% - Giving a talk on some subject.
% - The talk is between 15min and 45min long.
% - Style is ornate.

% MODIFIED by Jonathan Kew, 2008-07-06
% The header comments and encoding in this file were modified for inclusion with TeXworks.
% The content is otherwise unchanged from the original distributed with the beamer package.

\documentclass{beamer}
\setbeamercovered{invisible}



% Copyright 2004 by Till Tantau <tantau@users.sourceforge.net>.
%
% In principle, this file can be redistributed and/or modified under
% the terms of the GNU Public License, version 2.
%
% However, this file is supposed to be a template to be modified
% for your own needs. For this reason, if you use this file as a
% template and not specifically distribute it as part of a another
% package/program, I grant the extra permission to freely copy and
% modify this file as you see fit and even to delete this copyright
% notice. 


\mode<presentation>
{
  \usetheme{Warsaw}
  % or ...

  % or whatever (possibly just delete it)
}


\usepackage[english]{babel}
% or whatever

\usepackage[utf8]{inputenc}
\usepackage{qrcode}
\usepackage{hyperref}
\usepackage{ytableau}
\usepackage{longtable}
% or whatever

\usepackage{times}
\usepackage[T1]{fontenc}
% Or whatever. Note that the encoding and the font should match. If T1
% does not look nice, try deleting the line with the fontenc.

%%% MATH RELATED

\title{MATH 180 Strategies of Problem Solving} 
\subtitle
{} % (optional)

\author[W.R. Casper] % (optional, use only with lots of authors)
{W.R. Casper}
% - Use the \inst{?} command only if the authors have different
%   affiliation.

\institute[California State University Fullerton] % (optional, but mostly needed)
{
  Department of Mathematics\\
  California State University Fullerton}
% - Use the \inst command only if there are several affiliations.
% - Keep it simple, no one is interested in your street address.

\subject{Talks}
% This is only inserted into the PDF information catalog. Can be left
% out. 



% If you have a file called "university-logo-filename.xxx", where xxx
% is a graphic format that can be processed by latex or pdflatex,
% resp., then you can add a logo as follows:

% \pgfdeclareimage[height=0.5cm]{university-logo}{university-logo-filename}
% \logo{\pgfuseimage{university-logo}}



% Delete this, if you do not want the table of contents to pop up at
% the beginning of each subsection:
\AtBeginSubsection[]
{
  \begin{frame}<beamer>{Outline}
    \tableofcontents[currentsection,currentsubsection]
  \end{frame}
}


% If you wish to uncover everything in a step-wise fashion, uncomment
% the following command: 

%\beamerdefaultoverlayspecification{<+->}


\begin{document}

\begin{frame}
  \titlepage
\end{frame}

\begin{frame}{Outline}
  \tableofcontents
  % You might wish to add the option [pausesections]
\end{frame}

% Since this a solution template for a generic talk, very little can
% be said about how it should be structured. However, the talk length
% of between 15min and 45min and the theme suggest that you stick to
% the following rules:  

% - Exactly two or three sections (other than the summary).
% - At *most* three subsections per section.
% - Talk about 30s to 2min per frame. So there should be between about
%   15 and 30 frames, all told.

\section{SoPS Lecture 4}
\subsection{$3\times n$ Chomp}


\begin{frame}{Chomp Review}
  \begin{block}{Last time}
    \begin{itemize}
\pause
      \item we learned the "strategy-stealing argument"
\pause
      \item we proved {\color{blue}Player 1} has a perfect strategy
\pause
      \item we figured out the strategy for $1\times n$ and $2\times n$ Chomp
\pause
      \item by \textbf{symmetry}: strategy for $n\times 1$ and $n\times 2$ Chomp
    \end{itemize}
  \end{block}
\pause
  \begin{block}{Question}
    What about $3\times n$ chomp?
  \end{block}
\end{frame}

\begin{frame}{Perfect Strategy}
\pause
  \begin{block}{Question}
    What is the perfect strategy for $3\times 3$ Chomp?
    \begin{center}
      \ydiagram{3,3,3}
    \end{center}
  \end{block}
\end{frame}

\begin{frame}{Perfect Strategy}
\pause
  \begin{block}{Answer}
    First move:
    \begin{center}
      \ydiagram{1,1,3}
    \end{center}
\pause
    Then Copycat Strategy!
  \end{block}
\end{frame}

\begin{frame}{Perfect Strategy}
  \begin{block}{Question}
    What is the best first move for $3\times 4$ Chomp?
    \begin{center}
      \ydiagram{4,4,4}
    \end{center}
  \end{block}
\pause
  \begin{block}{Activity (10 mins)}
    Break into small groups and discuss.  What's the strategy?
  \end{block}
\end{frame}

\begin{frame}{Perfect Strategy}
  \begin{block}{First move?}
      Can't be left column or bottom row, because it'd make a rectangle (winning board), eg.
      $$(1,1): \ydiagram{4,4,4}\rightarrow \ydiagram{0,4,4}$$
      $$(3,3): \ydiagram{4,4,4}\rightarrow \ydiagram{2,2,2}$$
  \end{block}
\end{frame}

\begin{frame}{Perfect Strategy}
  \begin{block}{First move?}
      Can't be $(2,2)$, since then they could follow up with $(3,4)$
      $$(2,2): \ydiagram{4,4,4}\rightarrow \ydiagram{1,1,4}$$
      $$(3,4): \ydiagram{1,1,4}\rightarrow \ydiagram{1,1,3}$$
      Then we lose to Copycat Strategy
  \end{block}
\end{frame}

\begin{frame}{Perfect Strategy}
  \begin{block}{First move?}
      Can't be $(1,2)$, since then they could follow up with $(3,3)$
      $$(1,2): \ydiagram{4,4,4}\rightarrow \ydiagram{1,4,4}$$
      $$(3,3): \ydiagram{1,4,4}\rightarrow \ydiagram{1,2,2}$$
      Then we lose to the $n\times 2$ strategy
  \end{block}
\end{frame}

\begin{frame}{Perfect Strategy}
  \begin{block}{Theorem}
    The best first move for $3\times 4$ Chomp is $(2,3)$.
  \end{block}
\pause
  \begin{block}{Proof}
    Board after the first move:
    $$\ydiagram{2,2,4}$$
\pause
    How could {\color{red}Player 2} follow up?
  \end{block}
\end{frame}

\begin{frame}{Perfect Strategy}
  \begin{block}{Cases}
    $$(1,1): \ydiagram{2,2,4} \rightarrow \ydiagram{0,2,4}$$
    $$(2,1): \ydiagram{2,2,4} \rightarrow \ydiagram{0,0,4}$$
  \end{block}
\end{frame}

\begin{frame}{Perfect Strategy}
  \begin{block}{Cases}
    $$(1,2): \ydiagram{2,2,4} \rightarrow \ydiagram{1,2,4}$$
    $$(2,2): \ydiagram{2,2,4} \rightarrow \ydiagram{1,1,4}$$
    $$(3,2): \ydiagram{2,2,4} \rightarrow \ydiagram{1,1,1}$$
  \end{block}
\end{frame}

\begin{frame}{Perfect Strategy}
  \begin{block}{Cases}
    $$(3,3): \ydiagram{2,2,4} \rightarrow \ydiagram{2,2,2}$$
    $$(3,4): \ydiagram{2,2,4} \rightarrow \ydiagram{2,2,3}$$
  \end{block}
\end{frame}

\begin{frame}{Perfect Strategy}
  \begin{block}{Question}
    What is the best first move for $3\times 5$ Chomp?
    \begin{center}
      \ydiagram{5,5,5}
    \end{center}
  \end{block}
  \begin{block}{Activity (20 mins)}
    Break into small groups and discuss.  How many moves can you eliminate?
  \end{block}
\end{frame}

\begin{frame}{Perfect Strategy}
  \begin{block}{Process of elimination}
    \begin{center}
      \begin{ytableau}
        x &   &   &   &   \\
        x &   &   &   &   \\
        x & x & x & x & x
      \end{ytableau}
    \end{center}
    Avoid left side and bottom due to rectangles.
  \end{block}
\end{frame}

\begin{frame}{Perfect Strategy}
  \begin{block}{What about $(2,2)$?}
    \begin{center}
      \ydiagram{1,1,5}
    \end{center}
\pause
    Then they could just do the Copycat Strategy.
  \end{block}
\end{frame}

\begin{frame}{Perfect Strategy}
  \begin{block}{Process of elimination}
    \begin{center}
      \begin{ytableau}
        x &   &   &   &   \\
        x & x &   &   &   \\
        x & x & x & x & x
      \end{ytableau}
    \end{center}
  \end{block}
\end{frame}

\begin{frame}{Perfect Strategy}
  \begin{block}{What about $(1,2)$?}
    \begin{center}
      \ydiagram{1,5,5}
    \end{center}
\pause
    They can follow up with a $(3,3)$
    \begin{center}
      \ydiagram{1,2,2}
    \end{center}
    and then use $n\times 2$ tactics
  \end{block}
\end{frame}

\begin{frame}{Perfect Strategy}
  \begin{block}{Process of elimination}
    \begin{center}
      \begin{ytableau}
        x & x &   &   &   \\
        x & x &   &   &   \\
        x & x & x & x & x
      \end{ytableau}
    \end{center}
  \end{block}
\end{frame}

\begin{frame}{Perfect Strategy}
  \begin{block}{What about $(2,3)$?}
    \begin{center}
      \ydiagram{2,2,5}
    \end{center}
\pause
    They can follow up with a $(3,5)$
    \begin{center}
      \ydiagram{2,2,4}
    \end{center}
    Looks like they did the winning move for $3\times 4$
  \end{block}
\end{frame}

\begin{frame}{Perfect Strategy}
  \begin{block}{Process of elimination}
    \begin{center}
      \begin{ytableau}
        x & x &   &   &   \\
        x & x & x &   &   \\
        x & x & x & x & x
      \end{ytableau}
    \end{center}
  \end{block}
\end{frame}

\begin{frame}{Perfect Strategy}
  \begin{block}{What about $(1,3)$?}
    \begin{center}
      \ydiagram{2,5,5}
    \end{center}
\pause
    They can follow up with a $(2,4)$
    \begin{center}
      \ydiagram{2,3,5}
    \end{center}
    This is a losing board! 
  \end{block}
\end{frame}

\begin{frame}{Prove it}
  \begin{block}{Activity (10 mins)}
    Prove that the board
    \begin{center}
      \ydiagram{2,3,5}
    \end{center}
    is a \textbf{losing board}, ie. the next player to play can be forced to lose.
  \end{block}
\end{frame}

\begin{frame}{Perfect Strategy}
  \begin{block}{Process of elimination}
    \begin{center}
      \begin{ytableau}
        x & x & x &   &   \\
        x & x & x &   &   \\
        x & x & x & x & x
      \end{ytableau}
    \end{center}
  \end{block}
\end{frame}

\begin{frame}{Perfect Strategy}
  \begin{block}{What about $(2,4)$?}
    \begin{center}
      \ydiagram{3,3,5}
    \end{center}
\pause
    They can follow up with a $(1,3)$
    \begin{center}
      \ydiagram{2,3,5}
    \end{center}
    We just saw this was a losing position.
  \end{block}
\end{frame}

\begin{frame}{Perfect Strategy}
  \begin{block}{Process of elimination}
    \begin{center}
      \begin{ytableau}
        x & x & x &   &   \\
        x & x & x & x &   \\
        x & x & x & x & x
      \end{ytableau}
    \end{center}
  \end{block}
\end{frame}

\begin{frame}{Perfect Strategy}
  \begin{block}{Process of elimination}
    \begin{center}
      \begin{ytableau}
        x & x & x & ? & ? \\
        x & x & x & x & ? \\
        x & x & x & x & x
      \end{ytableau}
    \end{center}
  \end{block}
\end{frame}

\begin{frame}{Perfect Strategy}
  \begin{block}{Theorem}
    The perfect first move for $3\times 5$ Chomp is $(1,4)$.
    \begin{center}
      \ydiagram{3,5,5}
    \end{center}
  \end{block}
\end{frame}

\subsection{Strong Law of Small Numbers}

\begin{frame}{Getting moves for higher $n$ values}
  \begin{block}{Finding the perfect move requires ...}
    \begin{itemize}
\pause
      \item remembering lots of past boards
\pause
      \item checking lots of ``if then" cases
    \end{itemize}
  \end{block}
\pause
  \begin{block}{}
    \begin{center}
      Computers are really good at this!
    \end{center}
  \end{block}
\end{frame}

\begin{frame}{Computer calculating of best $3\times n$ move}
\begin{center}
\begin{tabular}{r|c|c}
\hline
$n$ & Width $w$ & Play at (row, col) \\
\hline


2  & 1  & (1,2)  \\
3  & 2  & (2,2)  \\
4  & 2  & (2,3)  \\
5  & 2  & (1,4)  \\
6  & 3  & (2,4)  \\
7  & 3  & (1,5)  \\
8  & 4  & (2,5)  \\
9  & 3  & (1,7)  \\
10 & 5  & (2,6)  \\
11 & 5  & (2,7)  \\
12 & 4  & (1,9)  \\
13 & 6  & (2,8)  \\
14 & 4  & (1,11) \\
15 & 7  & (2,9)  \\
\end{tabular}
\end{center}
\end{frame}

\begin{frame}{Positions of Best First Moves}
\begin{center}
\begin{tabular}{|c|c|c|c|c|c|c|c|c|c|c|}
\hline
    & 2 &   & 5 & 7 &   & 9 &   &12 &   &14 \\\hline
  1 & 3 & 4 & 6 & 8 &10 &11 &13 &15 &16 &18 \\\hline
    &   &   &   &   &   &   &   &   &   &   \\\hline
\end{tabular}
\end{center}
Entries represent the board length for which the corresponding spot is the best first move.
\begin{block}{Question}
Can you see any patterns?
\end{block}
\end{frame}

\begin{frame}{Notice patterns}
  \begin{block}{Submit your findings:}
    \begin{center}
      \qrcode[hyperlink,height=2cm]{https://padlet.com/williamrileycasper/pinching-pennies-losing-positions-h5xyswjj3ku3iqof}
    \end{center}
  \end{block}
\end{frame}

\begin{frame}{Create conjectures}
  \begin{block}{Conjecture 1}
    Each intege $n$ end up in a unique spot in the table.
  \end{block}
\pause
  \begin{block}{Conjecture 2}
    The numbers in any row get bigger from left to right
  \end{block}
\pause
  \begin{block}{Conjecture 3}
    The second row doesn't have any ``gaps"
  \end{block}
\pause
  \begin{block}{Next steps}
    \begin{itemize}
\pause
      \item gather more evidence
\pause
      \item try to prove it!
    \end{itemize}
  \end{block}
\end{frame}

\begin{frame}{Computer calculating of best $3\times n$ move}
\begin{center}
\begin{tabular}{r|c}
\hline
$n$ & Play at (row, col) \\
\hline
2  & ( 1 , 2 ) \\
3  & ( 2 , 2 ) \\
4  & ( 2 , 3 ) \\
5  & ( 1 , 4 ) \\
6  & ( 2 , 4 ) \\
7  & ( 1 , 5 ) \\
8  & ( 2 , 5 ) \\
9  & ( 1 , 7 ) \\
10  & ( 2 , 6 ) \\
11  & ( 2 , 7 ) \\
12  & ( 1 , 9 ) \\
13  & ( 2 , 8 ) \\
14  & ( 1 , 11 ) \\
\end{tabular}
\end{center}
\end{frame}
\begin{frame}{Computer calculating of best $3\times n$ move}
\begin{center}
\begin{tabular}{r|c}
\hline
$n$ & Play at (row, col) \\
\hline
15  & ( 2 , 9 ) \\
16  & ( 2 , 10 ) \\
17  & ( 1 , 13 ) \\
18  & ( 2 , 11 ) \\
19  & ( 1 , 14 ) \\
20  & ( 2 , 12 ) \\
21  & ( 2 , 13 ) \\
22  & ( 1 , 16 ) \\
23  & ( 1 , 17 ) \\
24  & ( 2 , 14 ) \\
25  & ( 2 , 15 ) \\
26  & ( 1 , 19 ) \\
27  & ( 2 , 16 ) \\
\end{tabular}
\end{center}
\end{frame}
\begin{frame}{Computer calculating of best $3\times n$ move}
\begin{center}
\begin{tabular}{r|c}
\hline
$n$ & Play at (row, col) \\
\hline
28  & ( 2 , 17 ) \\
29  & ( 1 , 21 ) \\
30  & ( 2 , 18 ) \\
31  & ( 1 , 22 ) \\
32  & ( 2 , 19 ) \\
33  & ( 1 , 24 ) \\
34  & ( 2 , 20 ) \\
35  & ( 2 , 21 ) \\
36  & ( 1 , 26 ) \\
37  & ( 2 , 22 ) \\
38  & ( 1 , 28 ) \\
39  & ( 2 , 23 ) \\
40  & ( 2 , 24 ) \\
\end{tabular}
\end{center}
\end{frame}
\begin{frame}{Computer calculating of best $3\times n$ move}
\begin{center}
\begin{tabular}{r|c}
\hline
$n$ & Play at (row, col) \\
\hline
41  & ( 1 , 30 ) \\
42  & ( 2 , 25 ) \\
43  & ( 1 , 31 ) \\
44  & ( 2 , 26 ) \\
45  & ( 2 , 27 ) \\
46  & ( 1 , 33 ) \\
47  & ( 1 , 34 ) \\
48  & ( 2 , 28 ) \\
49  & ( 2 , 29 ) \\
50  & ( 1 , 36 ) \\
51  & ( 2 , 30 ) \\
52  & ( 1 , 38 ) \\
\end{tabular}
\end{center}
\end{frame}
\begin{frame}{Computer calculating of best $3\times n$ move}
\begin{center}
\begin{tabular}{r|c}
\hline
$n$ & Play at (row, col) \\
\hline
53  & ( 2 , 31 ) \\
54  & ( 2 , 32 ) \\
55  & ( 1 , 40 ) \\
56  & ( 2 , 33 ) \\
57  & ( 2 , 34 ) \\
58  & ( 1 , 42 ) \\
59  & ( 2 , 35 ) \\
60  & ( 2 , 36 ) \\
61  & ( 1 , 44 ) \\
62  & ( 1 , 45 ) \\
63  & ( 2 , 37 ) \\
64  & ( 2 , 38 ) \\
65  & ( 1 , 47 ) \\
\end{tabular}
\end{center}
\end{frame}


\begin{frame}{Prove conjectures}
  \begin{block}{Theorem}
    Each intege $n$ end up in a unique spot in the table.
  \end{block}
\pause
  \begin{block}{Proof}
    Assume there is a spot $(j,k)$ which occurs twice, for values $m,n$ with $0<m<n$.\\
\pause
    Start with a $3\times n$ grid.\\
\pause
    Then if {\color{blue}Player 1} chooses $(j,k)$ they should win.\\
\pause
    However, if {\color{red}Player 2} follows by choosing $(3,n-m+1)$, it's equivalent to them having choosen $(j,k)$ on the $3\times n$ grid.\\
\pause
    Therefore {\color{red}Player 2} should be the one that wins.\\
\pause
    Contradiction.
  \end{block}
\end{frame}

\begin{frame}{Or find a counter-example}
\begin{center}
\begin{tabular}{r|c}
\hline
$n$ & Play at (row, col) \\
\hline
66  & ( 2 , 39 ) \\
67  & ( 1 , 48 ) \\
68  & ( 2 , 40 ) \\
69  & ( 1 , 50 ) \\
70  & ( 2 , 41 ) \\
71  & ( 1 , 51 ) \\
72  & ( 2 , 42 ) \\
73  & ( 2 , 43 ) \\
74  & ( 2 , 44 ) \\
75  & ( 1 , 53 ) \\
76  & ( 2 , 45 ) \\
77  & ( 1 , 55 ) \\
78  & ( 2 , 46 ) \\
\end{tabular}
\end{center}
\end{frame}
\begin{frame}{Computer calculating of best $3\times n$ move}
\begin{center}
\begin{tabular}{r|c}
\hline
$n$ & Play at (row, col) \\
\hline
79  & ( 1 , 56 ) \\
80  & ( 2 , 47 ) \\
81  & ( 1 , 58 ) \\
82  & ( 2 , 48 ) \\
83  & ( 2 , 49 ) \\
84  & ( 1 , 60 ) \\
85  & ( 2 , 50 ) \\
86  & ( 1 , 62 ) \\
87  & ( 2 , 51 ) \\
88  & ( 2 , 53 ) \\
89  & ( 2 , 52 ) \\
90  & ( 1 , 64 ) \\
91  & ( 1 , 65 ) \\
\end{tabular}
\end{center}
\end{frame}
\begin{frame}{Computer calculating of best $3\times n$ move}
\begin{center}
\begin{tabular}{r|c}
\hline
$n$ & Play at (row, col) \\
\hline
92  & ( 2 , 54 ) \\
93  & ( 2 , 55 ) \\
94  & ( 1 , 67 ) \\
95  & ( 2 , 56 ) \\
96  & ( 2 , 57 ) \\
97  & ( 1 , 69 ) \\
98  & ( 2 , 58 ) \\
99  & ( 1 , 71 ) \\
\end{tabular}
\end{center}
\end{frame}

\begin{frame}{Counter-example}
  \begin{block}{Conjecture}
    The number in the middle row must increase.
  \end{block}
\pause
  \begin{block}{First counter-example}
    For $n=88$, the best first move is $(2,53)$, but for $n=87$ the best first move is $(2,52)$
  \end{block}
\end{frame}

\begin{frame}{Strong Law of Small Numbers}
\pause
  \begin{block}{Richard Guy:}
    There aren't enough small numbers to make the many demands of them.
  \end{block}
\pause
  \begin{block}{Meaning}
    \begin{itemize}
\pause
      \item Mathematics is full of conjectures which are true for small numbers, but then fail for larger ones.
\pause
      \item Observations are evidence, but not proofs!
\pause
      \item Be suspicious!
    \end{itemize}
  \end{block}
\end{frame}

\begin{frame}{Strong Law of Small Numbers}
  \begin{block}{Example (Primes)}
    \begin{itemize}
\pause
      \item $31$ is prime
\pause
      \item $331$ is prime
\pause
      \item $3331$ is prime
\pause
      \item $33331$ is prime
\pause
      \item $333331$ is prime
\pause
      \item $3333331$ is prime
\pause
      \item $3333331$ is prime
\pause
      \item $33333331$ is prime
    \end{itemize}
\pause
    $$333333331 = 17\times 19607843$$
  \end{block}
\end{frame}

\begin{frame}{Strong Law of Small Numbers}
  \begin{block}{Example (Diophantine Equations)}
    For a long time, it was conjectured that
    $$x^4 + y^4 + z^4 = w^4$$
    has no nonzero **integer-valued** solutions.
  \end{block}
\pause
  \begin{block}{First counter-example}
    $$95800^4+217519^4+414560^4=4224814^4$$
  \end{block}
\end{frame}

\begin{frame}{Strong Law of Small Numbers}
\pause
  \begin{block}{Example (Polya's Conjecture)}
\pause
    An integer like $50 = 2\times 5\times 5$ has an \textbf{odd number of prime factors}.\\
\pause
    An integer like $16 = 2\times 2\times 2\times 2$ has an \textbf{even number of prime factors}.\\
\pause
    \textbf{Polya:} for any integer $n$, at least half of the integers $1 < k < n$ have an odd number of prime factors.
  \end{block}
\pause
  \begin{block}{First counter-example}
    $$n=906150257$$
  \end{block}
\end{frame}

\begin{frame}{Strong Law of Small Numbers}
\pause
  \begin{block}{Example (Goldbach Conjecture)}
\pause
    $12 = 7 + 5$, $30 = 13+17$, $46 = 23 + 23$.
\pause
    \textbf{Goldbach:} any even integer greater than $2$ is the sum of two primes
  \end{block}
\pause
  \begin{block}{Checked up through:}
    $$n=4,000,000,000,000,000,000$$
  \end{block}
\pause
$$\text{\color{red}no known counter-examples}$$
\end{frame}

\begin{frame}{Strong Law of Small Numbers}
\pause
  \begin{block}{Example (Collatz Conjecture)}
\pause
    $f(x) = 3x+1$ if $x$ is odd
\pause
    $f(x) = x/2$ if $x$ is even
\pause 
    $$13\mapsto 40\mapsto 20\mapsto 10\mapsto 5\mapsto 16\mapsto 8\mapsto 4\mapsto 2\mapsto 1$$
\pause
    \textbf{Goldbach:} any positive integer eventually goes to $1$
  \end{block}
\pause
  \begin{block}{Checked up through:}
    $$n=295,147,905,179,352,825,856$$
  \end{block}
\pause
$$\text{\color{red}no known counter-examples}$$
\end{frame}

\subsection{Wrapping up}

\begin{frame}{Discussion}
\pause
  \begin{block}{Conjecture vs. Proof}
    In science evidence is enough.
    In mathematics, we require proof because evidence can mislead us.
  \end{block}
\pause
  \begin{block}{Question 1}
    Is the best first move on a $3\times n$ easy to predict?
  \end{block}
\pause
  \begin{block}{Question 2}
    What does the Strong Law of Small Numbers mean in your own words?
    Can you think of any examples of the Strong Law of Small Numbers in action?
  \end{block}
\end{frame}


\begin{frame}{Question of the Day}
\pause
  \begin{block}{Puzzle}
    Suppose you start with a $3\times 4$ board.
    How many possible board shapes can we make by doing only Chomp moves.
  \end{block}
\pause
  \begin{block}{Activity (20 mins)}
    Discuss the puzzle with your neighbors.  What is the solution?
  \end{block}
\end{frame}

\begin{frame}{Final thoughts}
  \begin{block}{Reflection}
    \begin{itemize}
\pause
      \item Looking at the Chomp data above, come up with your own Conjecure about the best first move on a $3\times n$ Chomp board.
\pause
      \item Describe the Strong Law of Small Numbers in your own terms.  What does it say about evidence in mathematics.  Does it mean that evidence is useless?
    \end{itemize}
  \end{block}
\begin{center}
Submit your responses by \textbf{Tonight!}
\end{center}
\end{frame}


\end{document}


